% Head types in elevation. Common shank diameter d = 1.
\begin{tikzpicture}[>=Latex,line join=round]
\tikzset{
  steel/.style={fill=blue!8,draw=blue!55!black,thick},
  lbl/.style={font=\scriptsize,align=center,text width=3.6cm},
  sub/.style={font=\tiny\itshape,align=center,text width=3.6cm,text=black!60},
  dimn/.style={<->,gray!80!black,thin},
  note/.style={font=\tiny,text=red!65!black,align=center,text width=2.9cm}
}

% \shank{x centre}{total length}{fraction of length that is plain shank}
\newcommand{\shank}[3]{%
  \draw[steel] (#1-0.5,0) rectangle (#1+0.5,-#2);
  \foreach \k in {0,1,...,9}{%
    \pgfmathsetmacro\yy{-#2*(#3+(1-#3)*(0.07+0.098*\k))}
    \ifdim\yy pt>-#2 pt
      \draw[blue!55!black,thin] (#1-0.5,\yy) -- (#1+0.5,\yy-0.085);
    \fi}
  \draw[blue!55!black,thin,dashed] (#1-0.5,-#2*#3) -- (#1+0.5,-#2*#3);
}

%==================== ROW 1 : protruding heads ====================
\begin{scope}

% --- 1. hex head, tension ---
\def\x{0}
\shank{\x}{2.00}{0.30}
\fill[blue!8,draw=blue!55!black,thick]
  (\x-0.78,0) -- (\x+0.78,0) -- (\x+0.78,0.56) -- (\x+0.62,0.70)
  -- (\x-0.62,0.70) -- (\x-0.78,0.56) -- cycle;
\draw[blue!55!black] (\x-0.26,0) -- (\x-0.26,0.70);
\draw[blue!55!black] (\x+0.26,0) -- (\x+0.26,0.70);
\node[lbl] at (\x,-2.70) {\textbf{hex head} --- tension};
\node[sub] at (\x,-3.35) {full head height; long thread.\\ The load path is along the bolt axis.};

% --- 2. hex head, shear ---
\def\x{3.9}
\shank{\x}{2.00}{0.72}
\fill[blue!8,draw=blue!55!black,thick]
  (\x-0.78,0) -- (\x+0.78,0) -- (\x+0.78,0.28) -- (\x+0.66,0.40)
  -- (\x-0.66,0.40) -- (\x-0.78,0.28) -- cycle;
\draw[blue!55!black] (\x-0.26,0) -- (\x-0.26,0.40);
\draw[blue!55!black] (\x+0.26,0) -- (\x+0.26,0.40);
\node[lbl] at (\x,-2.70) {\textbf{hex head} --- shear};
\node[sub] at (\x,-3.35) {reduced head, short thread, so the\\ \emph{plain} shank fills the hole and the\\ shear plane never falls in a thread.};

% --- 3. twelve-point external wrenching ---
\def\x{7.8}
\shank{\x}{2.00}{0.30}
\fill[blue!8,draw=blue!55!black,thick]
  (\x-0.72,0) -- (\x+0.72,0) -- (\x+0.72,0.80) -- (\x-0.72,0.80) -- cycle;
\foreach \k in {-0.48,-0.24,0,0.24,0.48}{
  \draw[blue!55!black,thin] (\x+\k,0.03) -- (\x+\k,0.80);}
\node[lbl] at (\x,-2.70) {\textbf{twelve-point}};
\node[sub] at (\x,-3.35) {external wrenching: high preload in\\ a small wrench swing, in a\\ congested bay.};

% --- 4. internal wrenching / socket cap ---
\def\x{11.7}
\shank{\x}{2.00}{0.30}
\fill[blue!8,draw=blue!55!black,thick]
  (\x-0.74,0) -- (\x+0.74,0) -- (\x+0.74,0.94) -- (\x-0.74,0.94) -- cycle;
\fill[white,draw=blue!55!black,thick]
  (\x-0.34,0.95) -- (\x-0.34,0.42) -- (\x+0.34,0.42) -- (\x+0.34,0.95) -- cycle;
\node[lbl] at (\x,-2.70) {\textbf{internal wrenching}};
\node[sub] at (\x,-3.35) {socket head: no wrench clearance\\ needed \emph{around} the head, only\\ above it.};

\end{scope}

%==================== ROW 2 : flush heads and the countersink ====================
\begin{scope}[yshift=-6.2cm]

% --- 5. 100 deg countersunk, correct ---
\def\x{0}
\fill[black!6,draw=black!55] (\x-1.80,0) rectangle (\x+1.80,-0.64);
\fill[black!12,draw=black!55] (\x-1.80,-0.64) rectangle (\x+1.80,-1.40);
\fill[white] (\x-1.00,0.01) -- (\x-0.5,-0.42) -- (\x-0.5,-1.41) -- (\x+0.5,-1.41)
   -- (\x+0.5,-0.42) -- (\x+1.00,0.01) -- cycle;
\fill[blue!8,draw=blue!55!black,thick]
  (\x-1.00,0) -- (\x+1.00,0) -- (\x+0.5,-0.42) -- (\x+0.5,-1.95)
  -- (\x-0.5,-1.95) -- (\x-0.5,-0.42) -- cycle;
\draw[black!55] (\x-1.80,0) -- (\x-1.00,0);  \draw[black!55] (\x+1.00,0) -- (\x+1.80,0);
\draw[black!55] (\x-1.80,-0.64) -- (\x-0.5,-0.64); \draw[black!55] (\x+0.5,-0.64) -- (\x+1.80,-0.64);
% included angle, measured at the cone apex (x,-0.84)
\draw[gray!65,thin,dashed] (\x,-0.84) -- (\x-1.62,0.52);
\draw[gray!65,thin,dashed] (\x,-0.84) -- (\x+1.62,0.52);
\draw[dimn] ([shift=(40:1.18)]\x,-0.84) arc (40:140:1.18);
\node[font=\scriptsize,fill=white,inner sep=1pt] at (\x,0.46) {$100^\circ$};
% csk depth vs sheet thickness, on the left
\draw[dimn] (\x-1.52,0) -- (\x-1.52,-0.42);
\draw[dimn] (\x-2.05,0) -- (\x-2.05,-0.64);
\node[font=\tiny,anchor=east,text=gray!60!black] at (\x-2.10,-0.32) {$t$};
\node[font=\tiny,anchor=east,text=gray!60!black,fill=white,inner sep=0.5pt] at (\x-1.56,-0.21) {$h$};
\node[lbl] at (\x,-3.00) {\textbf{countersunk} (flush)};
\node[sub] at (\x,-3.65) {$100^\circ$ in aerospace. Keep\\ $h\le\tfrac23 t$: some parallel hole\\ must survive beneath the cone.};

% --- 6. knife edge, wrong ---
\def\x{5.2}
\fill[black!6,draw=black!55] (\x-1.80,0) rectangle (\x+1.80,-0.40);
\fill[black!12,draw=black!55] (\x-1.80,-0.40) rectangle (\x+1.80,-1.40);
\fill[white] (\x-1.00,0.01) -- (\x-0.5,-0.42) -- (\x-0.5,-1.41) -- (\x+0.5,-1.41)
   -- (\x+0.5,-0.42) -- (\x+1.00,0.01) -- cycle;
\fill[blue!8,draw=blue!55!black,thick]
  (\x-1.00,0) -- (\x+1.00,0) -- (\x+0.5,-0.42) -- (\x+0.5,-1.95)
  -- (\x-0.5,-1.95) -- (\x-0.5,-0.42) -- cycle;
\draw[black!55] (\x-1.80,0) -- (\x-1.00,0);  \draw[black!55] (\x+1.00,0) -- (\x+1.80,0);
\draw[black!55] (\x-1.80,-0.40) -- (\x-0.976,-0.40);
\draw[black!55] (\x+0.976,-0.40) -- (\x+1.80,-0.40);
\draw[red!70!black,very thick] (\x-0.53,-0.40) circle (0.19);
\draw[red!70!black,very thick] (\x+0.53,-0.40) circle (0.19);
\draw[red!65!black,thin] (\x-0.53,-0.59) -- (\x-1.05,-1.90);
\draw[red!65!black,thin] (\x+0.53,-0.59) -- (\x+1.05,-1.90);
\node[note,anchor=north] at (\x,-2.06) {the cone has cut clean\\ through the top sheet};
\node[lbl] at (\x,-3.00) {\textbf{knife edge} --- wrong};
\node[sub] at (\x,-3.65) {$h>t$: no parallel bearing length\\ is left, and the feather edge is\\ a built-in crack starter.};

% --- 7. dimpled thin sheet ---
\def\x{10.4}
\draw[black!55,fill=black!6]
  (\x-1.80,0) -- (\x-1.00,0) -- (\x-0.5,-0.43) -- (\x+0.5,-0.43) -- (\x+1.00,0)
  -- (\x+1.80,0) -- (\x+1.80,-0.20) -- (\x+1.00,-0.20) -- (\x+0.5,-0.63)
  -- (\x-0.5,-0.63) -- (\x-1.00,-0.20) -- (\x-1.80,-0.20) -- cycle;
\fill[black!12,draw=black!55]
  (\x-1.80,-0.20) -- (\x-1.00,-0.20) -- (\x-0.5,-0.63) -- (\x-0.5,-1.40)
  -- (\x+0.5,-1.40) -- (\x+0.5,-0.63) -- (\x+1.00,-0.20) -- (\x+1.80,-0.20)
  -- (\x+1.80,-1.40) -- (\x-1.80,-1.40) -- cycle;
\fill[blue!8,draw=blue!55!black,thick]
  (\x-1.00,0) -- (\x+1.00,0) -- (\x+0.5,-0.43) -- (\x+0.5,-1.95)
  -- (\x-0.5,-1.95) -- (\x-0.5,-0.43) -- cycle;
\node[lbl] at (\x,-3.00) {\textbf{dimpled} --- thin sheet};
\node[sub] at (\x,-3.65) {the sheet is \emph{formed}, not cut:\\ full thickness is kept and the\\ two sheets nest into each other.};

\end{scope}

\end{tikzpicture}
